% Round 16. Independently reviewed extension of the point-slice route.
% No historical claim is changed.
\section{Positive-dimensional proper-intersection residues}
\label{r16hp:section}

The point-slice theorem admits the following extension. Its only extra
analytic calculation is a compact-test version of the frozen Chern-degree
lemma PC6. The complete mathematical source, including the precise
archimedean pullback input, passed independent model-branch review at
SHA256 \path{de3d54cc4c04fa3b11f9c8de9c2ac71706dd1323ae88974fea8e9af56308550a},
recorded in
\path{round16/audits/proper_intersection_height_model_review.md}.
The corollary below additionally uses the subsequently reviewed proper
kernel covariance theorem to deduce its two image-domain memberships.

\begin{lemma}[Proper Chern degree for quasiprojective varieties]
\label{r16hp:degree}
Let $h:B\to A$ be a proper morphism of integral quasiprojective
$K$-varieties, with $\dim A=a$ and $\dim B=a+t$, $t\geq0$. Let $P$
be a homogeneous polynomial of degree $t$ in first Chern forms of
algebraic line bundles on $B$ with PWS metrics. For every compact actual
$\eta\in B_c^{a,a}(A^{\rm an})$,
\[
 \int_B P\wedge h^*\eta
      =\int_{h_*(P_{\rm alg}\cap[B])}\eta.
 \tag{HP16.1}\label{r16hp:degree-eq}
\]
The assertion is linear in cycles with their scheme multiplicities.
\end{lemma}

\begin{proof}
A quasiprojective variety has an ample invertible sheaf. Stacks 01Q3,
condition (7), implies that each of the finitely many line bundles in
$P$ is a quotient of two globally generated line bundles: for a large
ample power $H$, both $H$ and $L\otimes H$ are globally generated.
Choose a PWS metric on $H$ and the product metric on $L\otimes H$.
Their quotient is the original given metric on $L$, so this replacement
is an actual Chern identity. Expand $P$ by multilinearity.

For a resulting monomial, choose global sections of its globally
generated line factors successively, avoiding the generic points of
all the finitely many current cycle components. A finite-dimensional
space of sections generates at those points; over the infinite field
$K$, a finite union of proper linear subspaces cannot exhaust it.
Every divisor intersection is therefore proper. At each step use
scalar strong Poincare--Lelong with the remaining closed Chern factors
times $h^*\eta$. This test is compact, since $h$ is proper, and is
separately closed, since $\eta$ is top degree on $A$. The Green
correction consequently vanishes. Normalize a component first if
necessary; normalization does not change its actual top integral,
and divisor coefficients retain the generic lengths. The same argument
applies recursively to the resulting integral cycles.

The terminal cycle has dimension $a$. A component with lower-dimensional
image contributes zero, because the actual top test is zero on an
ordinary neighborhood of that algebraic image. On a dominating
component remove a proper algebraic closed subset of $A$ so that the
map is finite flat between smooth dense loci. Top-degree support
avoidance does not change either integral. The native finite degree
formula then gives its algebraic pushforward coefficient. Summing the
terminal components proves \eqref{r16hp:degree-eq}. Properness, rather
than projectivity over $K$, is the property needed for compact tests;
global generation is the only replacement required in the PC6 proof.
\end{proof}

\begin{theorem}[Curved residue under a proper intersection]
\label{r16hp:residue}
Use the supported-model-complex data $(G_C,\Omega_C,z_C)$ of
Theorem~\ref{r16h:curved-slice}, where $X$ is smooth projective
geometrically integral and $Z$ has pure codimension $r$. Let
$a:M\to X$ be a morphism from a normal integral quasiprojective
$K$-variety of dimension $m\geq r$. Suppose
\[
                   S=a^{-1}Z,\qquad \dim S\leq m-r.
\]
Then the ordinary actual form $a^*G_C$ on $M\setminus S$ satisfies
\[
 \int_{M\setminus S}a^*G_C\wedge D\eta
       +\int_{I_a(z_C)}\eta
       =\int_M a^*\Omega_C\wedge\eta,
 \qquad \eta\in B_c^{m-r,m-r}(M^{\rm an}),
 \tag{HP16.2}\label{r16hp:residue-eq}
\]
where $I_a(z_C)=\Delta_X^!([M]\times z_C)$ is the specified
top-dimensional cycle in $\operatorname{CH}_{m-r}(S)_{\mathbf Q}$.
For a global actual primitive $D\Psi=\Omega_C$, the form
$a^*(G_C-\Psi)$ has strong zero-curvature residue $-\delta_{I_a(z_C)}$
on all actual tests, with polar envelope $S$.
\end{theorem}

\begin{proof}
Use the same proper generically finite $f:T\to X$, the same supported
line-complex decomposition, and the same normalized finite trace on
$X\setminus Z$ as in Theorem~\ref{r16h:curved-slice}. Take the weighted
normalizations $N_j$ of the components of $T\times_X M$ dominating
$M$, with proper maps $b_j:N_j\to M$ and $\rho_j:N_j\to T$.
They are finite over $M\setminus S$, have dimension $m$, and
$\sum_j e_j(b_j)_*[N_j]=d[M]$. The pulled line maps have Cartier
divisors $D_{\alpha j}$ over $S$.

Scalar strong Poincare--Lelong on $N_j$ now uses the compact test
$\rho_j^*B_{r-1}(l_\alpha,m_\alpha)b_j^*\eta$, whose degree is
$(m-1,m-1)$. It gives exactly H16.6 with this new test degree.
The support of $D\eta$, of degree $(m-r+1,m-r+1)$, avoids $S$.
The finite-flat form base-change lemma H16.1 therefore identifies
the first term with $d\int a^*G_C D\eta$. The right term is the
pullback of an actual top form on $M$, so the native finite degree
formula identifies it with $d\int a^*\Omega_C\eta$.

For the residue terms, a divisor component has dimension $m-1$.
If its image in $S$ has dimension below $m-r$, the test vanishes
on an ordinary neighborhood of that image and its contribution is
zero. Otherwise its image is an integral component $A\subset S$
of dimension $m-r$. The map from that divisor component to $A$
is proper, and its Chern polynomial has degree $r-1$.
Lemma~\ref{r16hp:degree} identifies the analytic residue integral
with its algebraic Chern pushforward paired with the restricted
top actual test on $A$.

The supported Cartier-character identity, arbitrary localized base
change and the proper bivariant projection formula give the sum
of those algebraic cycles as
\[
 d\,\operatorname{ch}_r^S(La^*C)\cap[M]
                  =d\,\Delta_X^!([M]\times z_C).
\]
The last equality is the same regular-graph proof as H16.3; its
proof does not require $\dim M=r$. Because $\dim S\leq m-r$,
$\operatorname{CH}_{m-r}(S)$ is freely generated by the components
of that dimension. Thus this identifies the actual cycle
coefficients, not merely their image after forgetting support.
Divide by $d$, and then subtract the pulled global primitive.
\end{proof}

\paragraph{Support qualification.}
The displayed form is defined on exactly $M\setminus a^{-1}|Z|$.
The refined cycle can have vanishing coefficients or smaller support.
No removability across such canceled components is asserted. The
conclusion is deliberately a strong Green presentation with polar
envelope $a^{-1}|Z|$, which is precisely what disjoint-polar Stokes
needs when the other cycle avoids this set. It is not an exact-support
strong $G$ theorem for every possible refined pullback cycle.

\begin{corollary}[Projection formula on the proved height domains]
\label{r16hp:height-projection}
Let $f:M\to X$ be a morphism of smooth projective geometrically
integral varieties over a number field, of dimensions $m,n$.
Let $Z$ be a codimension-$a$ cycle on $M$, $W$ a codimension-$q$
cycle on $X$, with $a+q=m+1$, and suppose
\[
 |Z|\cap f^{-1}|W|=\varnothing,\qquad
 \dim f^{-1}|W|\leq m-q.
\]
Put $p=a+n-m$ and require $1\leq p\leq n$ and $1\leq q\leq m$.
Assume $[Z]\in\mathcal H^a(M)$ and $[W]\in\mathcal H^q(X)$,
for the domains of Theorem~\ref{r16h:global-height}. Then
\[
             h_X^{\rm an}([f_*Z],[W])
                       =h_M^{\rm an}([Z],[f^!W]).
 \tag{HP16.3}\label{r16hp:height-projection-eq}
\]
Here $f^!W$ denotes the specified proper refined-intersection cycle.
\end{corollary}

\begin{proof}
The new proper-kernel covariance theorem
\texttt{r16cov:proper-push} implies
$[f_*Z]\in\mathcal H^p(X)$ place by place. Frozen
contravariant functoriality of the obstruction, together with
$[f^!W]=f^*[W]$ in the Chow ring, implies
$[f^!W]\in\mathcal H^q(M)$. Geometric homological triviality
is preserved by these algebraic operations. Thus all four
classes needed for the two heights lie in their specified domains.

At each finite place choose the concrete Green form for $W$ from
the obstruction-zero construction. The preceding theorem makes its
pullback a strong Green presentation for $f^!W$ with envelope
$f^{-1}|W|$, disjoint from $|Z|$. Disjoint-polar Stokes compares
it with any chosen exact-support Green form for $f^!W$. The
ordinary cycle integration projection identity, using the already
proved local symmetry, gives
\[
 \lambda_v(f_*Z,W)=\int_{f_*Z}\gamma_W
                  =\int_Z f^*\gamma_W
                  =\lambda_v(Z,f^!W).
\]
For the middle equality, a component with lower-dimensional image
contributes zero by actual degree. A component with equal-dimensional
image is generically finite; top-degree support avoidance and the
native finite degree formula prove the identity with its full
pushforward coefficient. This is also the $t=0$ case of
Lemma~\ref{r16hp:degree}.

At an infinite place one uses a logarithmic-type Green form, rather
than pulling back an arbitrary current. Gillet--Soul\'e,
\emph{Arithmetic intersection theory}, Theorem~1.3.5, supplies a
logarithmic-type Green form for $W$ with smooth curvature. Its
curvature is cohomologically zero; the compact K\"ahler
$dd^c$ lemma in their Theorem~1.2.1 supplies a smooth global
primitive. Subtracting that primitive preserves logarithmic type
and gives zero curvature. Their Theorem~2.1.4(ii) applies to the
map $f$ and this chosen form: its pullback has the proper refined
intersection cycle as residue, together with a possible excess
current supported on the components of $f^{-1}|W|$ of codimension
strictly less than $q$. Our dimension hypothesis excludes all such
components, so the excess current is zero. Thus its pullback is a
Green presentation with envelope $f^{-1}|W|$, and classical
disjoint-polar Stokes and cycle integration give the same local
equality. The harmless scalar normalization of $dd^c$ is chosen
throughout to make the principal Green function $-\log|f|$.
The cited pullback theorem is the precise archimedean input;
no pullback of an arbitrary Green current is asserted.

All sums are finite by the preceding height theorem, so summation
proves \eqref{r16hp:height-projection-eq}.
\end{proof}

This corollary uses the independently proved proper covariance of
the obstruction \emph{kernel}, not a form-valued pushforward along
an arbitrary proper map. The new proper-intersection residue statement is an
identity tested against all actual forms, obtained from the concrete
supported-complex construction.
